% META: {"id": "TEXP-GRA-EX-009", "chapitre": "TEXP-GRAPHES", "type_objet": "exercice", "capacites": ["TEXP-GRA-C4"], "capacites_codes": ["C4"], "methodes": ["M4"], "parcours": 2, "competences": ["calculer"], "duree_min": 15, "mode_creation": "ex_nihilo", "sources_inspiration": [], "fichier_tex": "chapitres/TEXP-GRAPHES/exercices/TEXP-GRA-EX-009.tex", "corrige_tex": "chapitres/TEXP-GRAPHES/corriges/TEXP-GRA-CO-009.tex", "status": "generated"}
% BEGIN-VERIFY
% from sympy import Matrix
% aretes = [(1, 2), (2, 3), (3, 4), (1, 4), (1, 3)]
% M = Matrix(4, 4, lambda i, j: 1 if (i + 1, j + 1) in aretes or (j + 1, i + 1) in aretes else 0)
% M2, M3 = M ** 2, M ** 3
% assert M2 == Matrix([[3, 1, 2, 1], [1, 2, 1, 2], [2, 1, 3, 1], [1, 2, 1, 2]])
% assert M2[0, 2] == 2
% assert M3 == Matrix([[4, 5, 5, 5], [5, 2, 5, 2], [5, 5, 4, 5], [5, 2, 5, 2]])
% assert M3[0, 1] == 5
% assert M2[0, 0] == sum(M.row(0)) == 3
% chemins = [c for c in [(1, x, 3) for x in range(1, 5)]
%            if M[c[0] - 1, c[1] - 1] and M[c[1] - 1, c[2] - 1]]
% assert len(chemins) == 2
% END-VERIFY
\begin{exercice}{TEXP-GRA-EX-009}{2}{15}
On reprend le graphe de sommets $1,2,3,4$ et d'arêtes $\{1,2\}$, $\{2,3\}$,
$\{3,4\}$, $\{1,4\}$, $\{1,3\}$, de matrice d'adjacence $M$.
\begin{enumerate}
  \item Calculer $M^2$ et interpréter son coefficient $(1\,;\,3)$.
  \item Calculer $M^3$ et donner le nombre de chemins de longueur $3$ reliant
        $1$ à $2$.
  \item Expliquer pourquoi le coefficient $(1\,;\,1)$ de $M^2$ vaut le degré
        du sommet $1$.
\end{enumerate}
\end{exercice}
